\beforesec
\section{Limitations, impact, and conclusion}
\label{sec:conclusion}
\postsec
\smallpara{Limitations.} While we expect our findings to be more broadly applicable to other domains such as natural language processing, our investigation here is limited to image classification.
Exploring fine-tuning for object detection and natural language processing are interesting directions for future work.
Moreover, although the interpolation parameter setting $\alpha{=}0.5$ provides good overall performance, we leave the question of finding the optimal $\alpha$ for specific target distributions to future work.


\smallpara{Impact.}
 \citet{radford2021learning} and \citet{brown2020language} extensively discuss the broader impact of large zero-shot models and identify potential causes of harm including
 model biases and potential malicious uses such as surveillance systems.
 WiSE-FT is a fine-tuning method that builds on such models, and thus may perpetuate their negative impact.
 
 \smallpara{Conclusion.}
WiSE-FT can substantially improve performance under distribution shift with minimal or no loss in accuracy on the target distribution compared to standard fine-tuning.
We view WiSE-FT as a first step towards more sophisticated fine-tuning schemes and anticipate that future work will continue to leverage the robustness of zero-shot models for building more reliable neural networks.







\subsection*{Acknowledgements}
We thank
Anders Andreassen,
Tim Dettmers,
Jesse Dodge,
Katie Everett,
Samir Gadre,
Ari Holtzman,
Sewon Min,
Mohammad Norouzi,
Nam Pho, 
Ben Poole,
Sarah Pratt,
Alec Radford,
Jon Shlens,
and Rohan Taori for helpful discussions and draft feedback, Hyak at UW for computing support, 
Rosanne Liu for fostering the collaboration,
and Basil Mustafa for providing an earlier version of the mapping between JFT and ImageNet classes.
This work is in part supported by NSF IIS 1652052, IIS 17303166, DARPA N66001-19-2-4031, DARPA W911NF-15-1-0543 and gifts from Allen Institute for Artificial Intelligence.